\section{Класс с ограниченной начальной невязкой}\label{app:gap}

В основном тексте о начальной точке известно только расстояние $R$ до точки минимума, и начальная невязка по функции может быть экспоненциально большой по $L_1R$. В этом приложении рассмотрен класс $\mathcal P_\Delta$ из раздела~\ref{sec:setting}, в котором начальная невязка ограничена известным числом $\Delta$, и прежде всего случай $\Delta=4\varepsilon$. Показатель степени при $L_1R$ в оценке сложности падает при этом с единицы до $2/3$.

Условие $f(x^0)-f^\ast\le4\varepsilon$ естественно в схемах с рестартами: каждая фаза стартует из точки, где невязка уже ограничена постоянным кратным точности, требуемой от этой фазы.

Нижняя оценка доказана в подразделе~\ref{sec:eq-chain} на цепочке с равными наклонами звеньев. Верхние оценки получены вариантом ускоренного метода раздела~\ref{sec:upper} с другими способами подсчёта итераций (подразделы~\ref{sec:gap-method}--\ref{sec:count2}). Итоговые утверждения сформулированы в конце соответствующих подразделов.

\subsection{Нижняя оценка: цепочка с равными наклонами}\label{sec:eq-chain}

Как и в подразделе~\ref{sec:ordered}, построим цепочку сопряжением на многограннике и применим лемму~\ref{lem:conj}. Здесь наклоны всех звеньев равны, а сдвиги образуют арифметическую прогрессию с шагом $-b$.

\paragraph{Наклон и сдвиги.} Зафиксируем $n\ge1$ и положим
\begin{equation}\label{eq:eq-chain-params}
\begin{aligned}
 &b=R\Bigl(\sum_{k=n+1}^{2n}k^2\Bigr)^{-1/2},\\
 &\sigma_i=(2n-i)\,b,\\
 &G=\frac{2\varepsilon}{nb} .
\end{aligned}
\end{equation}
где $i=1,\dots,n$. Сумма в определении $b$ лежит между $n^3$ и $4n^3$.

\paragraph{Функция на многограннике.} Пусть функция $r\in C^1([0,1])$ сильно выпукла с некоторым положительным модулем, $r(0)=r'(0)=0$ и $r'(1)=b$. Рассмотрим многогранник
\[
 K=\bigl\{z\in\R^n:\ z\ge0,\ Z_n\le G\bigr\},
\]
где $Z_n=z_1+\dots+z_n$. Для $z\in K$ положим
\[
\begin{aligned}
 p_i&=\frac{z_i}G,\\
 p_0&=1-\frac{Z_n}G ,
\end{aligned}
\]
где $i=1,\dots,n$.
Координата $p_0$ --- остаток до ограничения $Z_n\le G$, она обеспечивает плоское множество точек минимума. Определим функцию на $K$ и сопряжённую к ней:
\begin{equation}\label{eq:eq-chain}
\begin{aligned}
 &H(z)=G\Bigl[\sum_{i=0}^nr(p_i)-r(1)\Bigr],\\
 &\psi(y)=\max_{z\in K}\bigl\{\ip{y+\sigma}z-H(z)\bigr\}.
\end{aligned}
\end{equation}
Здесь $H(0)=H(Ge_i)=0$.

\begin{lemma}[свойства цепочки с равными наклонами]\label{lem:eq-chain}
Функция $\psi$ выпукла, непрерывно дифференцируема, и
\begin{enumerate}
\item[(а)] её минимум равен нулю, а множество точек минимума равно
\[
 \bigl\{y\in\R^n:\ y_i\le-(\sigma_i+b)\ \text{при}\ i=1,\dots,n\bigr\}.
\]
Ближайшая к нулю точка минимума $-b\,(2n,2n-1,\dots,n+1)$ имеет норму $R$,
\item[(б)] $\psi(0)=(4-2/n)\varepsilon$,
\item[(в)] при $y_n=0$ имеем $\psi(y)\ge2\varepsilon$,
\item[(г)] $\psi$ --- цепочка,
\item[(д)] $\nabla\psi(y)=z(y)=G\,(p_1,\dots,p_n)$, где $z(y)$ --- единственная точка максимума в~\eqref{eq:eq-chain}.
\end{enumerate}
\end{lemma}

\begin{proof}
\leavevmode
\begin{itemize}
\item \emph{Градиент.} Сильная выпуклость $r$ даёт сильную выпуклость $H$ на $K$ с постоянным модулем. Поэтому точка максимума единственна, и, как в доказательстве леммы~\ref{lem:conj}, $\nabla\psi(y)=z(y)$ и градиент непрерывен. Это утверждение~(д).
\item \emph{Условия оптимальности.} Положим $\theta_0=0$ и $\theta_i=y_i+\sigma_i$ при $i\ge1$. Задача максимизации в~\eqref{eq:eq-chain} равносильна максимизации
\[
 \sum_{i=0}^n\theta_ip_i-\sum_{i=0}^nr(p_i)
\]
на симплексе. Перенос веса между двумя положительными координатами и из положительной координаты в нулевую даёт число $\nu$, для которого при $0\le i\le n$
\begin{equation}\label{eq:eq-kkt}
 \theta_i\ \begin{cases}=\nu+r'(p_i), & p_i>0,\\ \le\nu, & p_i=0.\end{cases}
\end{equation}
Эти условия также достаточны по выпуклости $r$. Так как $0\le r'(p_i)\le b$, всегда $\nu\ge\theta_i-b$.
\item \emph{Утверждение~(а).} Точка $z=0$ допустима, поэтому $\psi\ge0$. Равенство достигается ровно тогда, когда оптимален вектор $p=e_0$. По~\eqref{eq:eq-kkt} это равносильно $\nu=-b$ и $\theta_i\le-b$ при $i\ge1$. Формулы для множества точек минимума и ближайшей к нулю точки следуют из~\eqref{eq:eq-chain-params}.
\item \emph{Утверждение~(б).} При $y=0$ условия~\eqref{eq:eq-kkt} выполнены для $p=e_1$ и $\nu=\sigma_1-b=(2n-2)b$: все остальные $\theta_i$, включая $\theta_0=0$, не больше $\nu$. Поэтому $\psi(0)=G\sigma_1=(4-2/n)\varepsilon$.
\item \emph{Утверждение~(в).} Подставляя $z=Ge_n$, получаем $\psi(y)\ge G\sigma_n=2\varepsilon$.
\item \emph{Утверждение~(г).} Если $y_i=0$ при $i\ge j$, то при $l>j$
\[
 \theta_l=\sigma_l\le\sigma_j-b=\theta_j-b\le\nu .
\]
Положительное $p_l$ противоречило бы~\eqref{eq:eq-kkt}, поскольку $r'(p_l)>0$. Значит, $l$-я координата градиента $Gp_l$ равна нулю.\qedhere
\end{itemize}
\end{proof}

Для проверки условия класса достаточно одномерной оценки кривизны функции $r$.

\begin{lemma}[условие класса для равных наклонов]\label{lem:eq-conj-class}
Если $r$ дважды непрерывно дифференцируема на $(0,1)$ и при $0<p<1$
\begin{equation}\label{eq:eq-dual-curvature}
 r''(p)\ge\frac{G}{L_0+L_1Gp},
\end{equation}
то $\psi\in\F_{L_0,L_1}^n$.
\end{lemma}

\begin{proof}
Пусть $z,z'\in K$ различны и $u=z'-z$. На интервале $0<t<1$ координаты $p_i$ точки $z+tu$ либо положительны, либо тождественно равны нулю. Поэтому
\begin{align*}
 \frac{d^2}{dt^2}H(z+tu)
 &=\frac1G\sum_{i=1}^nr''(p_i)u_i^2+\frac1Gr''(p_0)U_n^2\\
 &\by{\ge}{\eqref{eq:eq-dual-curvature}}\frac{\norm u^2}{L_0+L_1\norm{z+tu}}
 \ge\frac{\norm u^2}{L_0+L_1\max\{\norm z,\norm{z'}\}} ,
\end{align*}
где $U_n=u_1+\dots+u_n$. Слагаемые с тождественно нулевыми координатами пропущены. Если $p_0\equiv0$, то $U_n=0$ и последнее слагаемое также отсутствует. Во втором переходе использовано неравенство $Gp_i=(z+tu)_i\le\norm{z+tu}$, в третьем --- выпуклость нормы. Интегрирование и непрерывность на концах дают условие леммы~\ref{lem:conj}, которая завершает доказательство.
\end{proof}

\paragraph{Квадратичная функция $r$.} Возьмём $r(p)=bp^2/2$. Тогда $r''=b$. Пусть
\begin{equation}\label{eq:eq-quad-cond}
 4n\le\sqrt{\frac{L_0R^2}\varepsilon} .
\end{equation}
Тогда
\begin{align*}
 \frac Gb
 &\by{=}{\eqref{eq:eq-chain-params}}\frac{2\varepsilon}{nR^2}\sum_{k=n+1}^{2n}k^2
 \le\frac{8\varepsilon n^2}{R^2}\\
 &\by{\le}{\eqref{eq:eq-quad-cond}}\frac{L_0}2 .
\end{align*}
Следовательно, условие~\eqref{eq:eq-dual-curvature} выполнено, и градиент функции $\psi$ даже липшицев с константой не больше $L_0/2$.

\paragraph{Логарифмическая функция $r$.} Положим
\begin{equation}\label{eq:eq-log-reg}
\begin{aligned}
 r(p)&=\frac1{L_1}\Bigl[(p+\chi)\log\Bigl(1+\frac p\chi\Bigr)-p\Bigr],\\
 \chi&=\frac1{e^{L_1b}-1} .
\end{aligned}
\end{equation}
Тогда
\[
\begin{aligned}
 &r'(p)=\frac1{L_1}\log\Bigl(1+\frac p\chi\Bigr),\\
 &r''(p)=\frac1{L_1(p+\chi)},
\end{aligned}
\]
и $r'(1)=b$. Условие~\eqref{eq:eq-dual-curvature} выполнено, если $GL_1\chi\le L_0$, поскольку тогда $GL_1(p+\chi)\le L_0+L_1Gp$. Проверим это при условиях
\begin{equation}\label{eq:eq-log-cond}
\begin{aligned}
 &L_1^2\varepsilon\ge L_0,\\
 &n\le\Bigl(\frac{L_1R}{8\log(2+L_1^2\varepsilon/L_0)}\Bigr)^{2/3} .
\end{aligned}
\end{equation}
Из~\eqref{eq:eq-chain-params} и~\eqref{eq:eq-log-cond}
\begin{align*}
 L_1b
 &\by{\ge}{\eqref{eq:eq-chain-params}}\frac{L_1R}{2n^{3/2}}
 \by{\ge}{\eqref{eq:eq-log-cond}}4\log\Bigl(2+\frac{L_1^2\varepsilon}{L_0}\Bigr)\\
 &\bypad{\ge}{\eqref{eq:eq-log-cond}}{\eqref{eq:eq-chain-params}}4\log3>4 .
\end{align*}
В последнем переходе использовано первое условие~\eqref{eq:eq-log-cond}.
Поэтому $e^{L_1b}\ge2$ и
\[
 \chi\le2e^{-L_1b}\le2\Bigl(2+\frac{L_1^2\varepsilon}{L_0}\Bigr)^{-2},
\]
и
\begin{align*}
 \frac{GL_1\chi}{L_0}
 &=\frac{2L_1^2\varepsilon}{nL_1b\,L_0}\,\chi
 \le\frac4{nL_1b}\cdot\frac{L_1^2\varepsilon/L_0}{\bigl(2+L_1^2\varepsilon/L_0\bigr)^2}\\
 &\by{\le}{(\mathrm a)}\frac1{2nL_1b}<\frac12 .
\end{align*}
Здесь (а) --- неравенство $t/(2+t)^2\le1/8$ при $t\ge0$, равносильное $(t-2)^2\ge0$.

\begin{theorem}[нижняя оценка на классе с ограниченной начальной невязкой]\label{thm:lower-gap}
Пусть натуральное число $n$ удовлетворяет условию~\eqref{eq:eq-quad-cond} или условию~\eqref{eq:eq-log-cond}. Тогда для каждого детерминированного метода с $n-1$ обращениями к оракулу в $\R^{2n}$ найдётся функция $f\in\mathcal P_{4\varepsilon}$ на $\R^{2n}$, для которой $f^\ast=0$, расстояние от нуля до множества точек минимума равно $R$, $f(0)=(4-2/n)\varepsilon$ и $f(\hat x)\ge2\varepsilon$. В частности, $\mathcal N(\mathcal P_{4\varepsilon})\ge n$, и при $L_0R^2\ge\varepsilon$
\begin{equation}\label{eq:lower-gap}
 1+\mathcal N(\mathcal P_{4\varepsilon})\ge\frac1{12}\Bigl[1+\sqrt{\frac{L_0R^2}\varepsilon}+\Bigl(\frac{L_1R}{\log(2+L_1^2\varepsilon/L_0)}\Bigr)^{2/3}\Bigr].
\end{equation}
\end{theorem}

\begin{proof}
\leavevmode
\begin{itemize}
\item \emph{Наихудшие функции.} Если выполнено условие~\eqref{eq:eq-quad-cond}, возьмём квадратичную функцию $r$, иначе --- логарифмическую~\eqref{eq:eq-log-reg}. В обоих случаях $\psi\in\F_{L_0,L_1}^n$ по лемме~\ref{lem:eq-conj-class}, а лемма~\ref{lem:eq-chain} даёт остальные условия предложения~\ref{prop:reduction}(б) с $\mathrm{gap}=2\varepsilon$ и $\Delta=4\varepsilon$. По этому предложению для каждого метода найдётся матрица $V$ с ортонормированными столбцами, для которой функция $f(x)=\psi(V^\top x)$ лежит в $\mathcal P_{4\varepsilon}$ и $f(\hat x)\ge2\varepsilon$. Остальные свойства $f$ дают утверждения~(а) и~(б) леммы~\ref{lem:eq-chain}.
\item \emph{Случай $L_1^2\varepsilon\ge L_0$.} Для каждого из условий~\eqref{eq:eq-quad-cond} и~\eqref{eq:eq-log-cond} возьмём в качестве $n$ целую часть наибольшей допустимой правой части. Если она положительна, первая часть теоремы даёт $\mathcal N(\mathcal P_{4\varepsilon})\ge n$, а если равна нулю, это неравенство очевидно. Вместе с неравенством $1+\lfloor x\rfloor\ge x$ получаем
\begin{align*}
 1+\mathcal N(\mathcal P_{4\varepsilon})&\ge1,\\
 1+\mathcal N(\mathcal P_{4\varepsilon})&\ge\frac14\sqrt{\frac{L_0R^2}\varepsilon},\\
 1+\mathcal N(\mathcal P_{4\varepsilon})&\ge\Bigl(\frac{L_1R}{8\log(2+L_1^2\varepsilon/L_0)}\Bigr)^{2/3}\\
 &=\frac14\Bigl(\frac{L_1R}{\log(2+L_1^2\varepsilon/L_0)}\Bigr)^{2/3} .
\end{align*}
Складывая эти неравенства с множителями $1$, $4$ и $4$, получаем, что скобка в~\eqref{eq:lower-gap} не больше $9\bigl(1+\mathcal N(\mathcal P_{4\varepsilon})\bigr)$.
\item \emph{Случай $L_1^2\varepsilon<L_0$.} Тогда $(L_1R)^2<L_0R^2/\varepsilon$, и при $L_0R^2\ge\varepsilon$
\begin{align*}
 \Bigl(\frac{L_1R}{\log(2+L_1^2\varepsilon/L_0)}\Bigr)^{2/3}
 &\by{\le}{(\mathrm a)}(\log2)^{-2/3}(L_1R)^{2/3}
 \by{<}{(\mathrm б)}\frac32\Bigl(\frac{L_0R^2}\varepsilon\Bigr)^{1/3}\\
 &\by{\le}{(\mathrm в)}\frac32\sqrt{\frac{L_0R^2}\varepsilon} ,
\end{align*}
где (а) --- неравенство $\log(2+t)\ge\log2$, (б) --- неравенства $(\log2)^{-2/3}<3/2$ и $(L_1R)^2<L_0R^2/\varepsilon$, (в) --- условие $L_0R^2\ge\varepsilon$. Значит, скобка в~\eqref{eq:lower-gap} не больше
\begin{align*}
 1+\frac52\sqrt{\frac{L_0R^2}\varepsilon}
 &\le\bigl(1+\mathcal N(\mathcal P_{4\varepsilon})\bigr)+10\bigl(1+\mathcal N(\mathcal P_{4\varepsilon})\bigr)\\
 &=11\bigl(1+\mathcal N(\mathcal P_{4\varepsilon})\bigr).
\end{align*}
Здесь использованы первые два неравенства из предыдущего случая: они получены без условия на знак $L_1^2\varepsilon-L_0$.
\end{itemize}
В обоих случаях подходит константа $1/12$.
\end{proof}

Показатель $2/3$ возникает так. Сдвиги $\sigma_i=(2n-i)b$ образуют арифметическую прогрессию с шагом $-b$, поэтому $R^2$ имеет порядок $b^2n^3$. Условие класса для логарифмической функции $r$ требует, чтобы $L_1b$ было не меньше логарифма, как в условии~\eqref{eq:eq-log-cond}. Эти две оценки дают длину цепочки порядка $(L_1R)^{2/3}$ с точностью до логарифма. Начальная невязка меньше $4\varepsilon$, и теорема~\ref{thm:upper-gap}(а) ниже показывает, что показатель $2/3$ точен с точностью до логарифмических множителей. Начальная невязка функций теоремы~\ref{thm:ordered}, напротив, не ограничена равномерно величиной $4\varepsilon$ (см. замечание после неё).


\subsection{Вариант метода без поправки}\label{sec:gap-method}

В этом подразделе и далее $f\in\F_{L_0,L_1}^d$, а $x^\ast$ --- точка минимума функции $f$. Используется вариант алгоритма~\ref{alg:am} без добавки $\Phi$, и с начальным весом $a_0\ge0$ (алгоритм~\ref{alg:am-gap}). Положим $A_k=a_0+a_1+\dots+a_k$. Коэффициент $a_k>0$ определяется уравнением
\begin{equation}\label{eq:coef-gap-rule}
 32\bigl(L_0+L_1\norm{\nabla f(y^k)}\bigr)a_k^2=A_{k-1}+a_k ,
\end{equation}
а число шагов одномерного поиска --- условием~\eqref{eq:search-m}.

\begin{algorithm}[ht]
\caption{Вариант метода SHERPA без поправки}\label{alg:am-gap}
\begin{algorithmic}[1]
\Require начальная точка $x^0$, начальный вес $a_0\ge0$, числа $r,\delta>0$, натуральное число $m$ с условием~\eqref{eq:search-m}, число итераций $K$
\State $v^0=x^0$, $A_0=a_0$
\For{$k=1,\dots,K$}
 \State $y^k$ --- результат алгоритма~\ref{alg:search} на отрезке $[v^{k-1},x^{k-1}]$ с числами $r$, $\delta$, $m$
 \If{сработал критерий остановки}
  \State \Return $y^k$
 \EndIf
 \State $a_k$ --- положительный корень уравнения~\eqref{eq:coef-gap-rule}, $A_k=A_{k-1}+a_k$
 \State $x^k=y^k-\dfrac{\nabla f(y^k)}{4(L_0+L_1\norm{\nabla f(y^k)})}$
 \State $v^k=v^{k-1}-a_k\nabla f(y^k)$
\EndFor
\State \Return $x^K$
\end{algorithmic}
\end{algorithm}

\begin{lemma}\label{lem:energy-gap}
Пусть $\norm{x^0-x^\ast}\le R_0$, начальная невязка $f(x^0)-f^\ast$ не больше $\Delta_0$ и
\[
 r^2\ge R_0^2+2a_0\Delta_0 .
\]
Если первые $k$ итераций алгоритма~\ref{alg:am-gap} прошли без срабатывания критерия остановки, то
\begin{equation}\label{eq:energy-gap}
 A_k\bigl(f(x^k)-f^\ast\bigr)+\frac12\sum_{i=1}^ka_i^2\norm{\nabla f(y^i)}^2+\frac12\norm{v^k-x^\ast}^2\le\frac12R_0^2+a_0\Delta_0 ,
\end{equation}
и все точки метода и одномерного поиска лежат в шаре $B(x^\ast,r)$. Кроме того, при всех $i\le k$
\begin{equation}\label{eq:gap-monotone}
\begin{aligned}
 f(x^{i-1})&\ge f(y^i),\\
 f(y^i)&\ge f(x^i)+\frac{\norm{\nabla f(y^i)}^2}{8\bigl(L_0+L_1\norm{\nabla f(y^i)}\bigr)} .
\end{aligned}
\end{equation}
\end{lemma}

\begin{proof}
Доказательство повторяет доказательство леммы~\ref{lem:energy}. При $k=0$ неравенство~\eqref{eq:energy-gap} следует из условий на $x^0$. Шаг индукции отличается в трёх местах.
\begin{itemize}
\item \emph{Замены.} Вместо леммы~\ref{lem:anchored} используется неравенство выпуклости $\ip{\nabla f(y^k)}{y^k-x^\ast}\ge f(y^k)-f^\ast$, а лемма~\ref{lem:search} используется так же. Неравенства~\eqref{eq:gap-monotone} следуют из неё и леммы~\ref{lem:safe-step}.
\item \emph{Приращение.} С этими заменами оценки расстояния и невязки из доказательства леммы~\ref{lem:energy} дают
\begin{align*}
 \MoveEqLeft\Bigl[A_k\bigl(f(x^k)-f^\ast\bigr)+\frac12\norm{v^k-x^\ast}^2\Bigr]-\Bigl[A_{k-1}\bigl(f(x^{k-1})-f^\ast\bigr)+\frac12\norm{v^{k-1}-x^\ast}^2\Bigr]\\
 &\bypad{\le}{}{\eqref{eq:coef-gap-rule}}\Bigl(\frac{a_k}{16}-\frac{5A_k}{24}\Bigr)\frac{\norm{\nabla f(y^k)}^2}{L_0+L_1\norm{\nabla f(y^k)}}+\frac12a_k^2\norm{\nabla f(y^k)}^2\\
 &\by{=}{\eqref{eq:coef-gap-rule}}\Bigl(\frac{a_k}{16}-\frac{5A_k}{24}+\frac{A_k}{32}\Bigr)\frac{\norm{\nabla f(y^k)}^2}{L_0+L_1\norm{\nabla f(y^k)}}-\frac12a_k^2\norm{\nabla f(y^k)}^2\\
 &\bypad{\le}{(\mathrm a)}{\eqref{eq:coef-gap-rule}}-\frac12a_k^2\norm{\nabla f(y^k)}^2 ,
\end{align*}
где (а) --- неравенства $a_k\le A_k$ и $1/16-5/24+1/32<0$.
\item \emph{Шар радиуса $r$.} По неравенству~\eqref{eq:energy-gap} и условию леммы
\[
 \norm{v^k-x^\ast}^2\le R_0^2+2a_0\Delta_0\le r^2 ,
\]
а точка $x^k$ лежит в шаре по лемме~\ref{lem:safe-step}.\qedhere
\end{itemize}
\end{proof}

Из уравнения~\eqref{eq:coef-gap-rule} и равенства $A_k=A_{k-1}+a_k$ при каждом $i\ge1$ следуют соотношения
\begin{equation}\label{eq:coef-gap}
\begin{aligned}
 a_i^2&=\frac{A_i}{32\bigl(L_0+L_1\norm{\nabla f(y^i)}\bigr)},\\
 \sqrt{A_i}-\sqrt{A_{i-1}}&=\frac{a_i}{\sqrt{A_i}+\sqrt{A_{i-1}}}\by{\ge}{A_{i-1}\le A_i}\frac{a_i}{2\sqrt{A_i}}\\
 &=\frac1{2\sqrt{32\bigl(L_0+L_1\norm{\nabla f(y^i)}\bigr)}} .
\end{aligned}
\end{equation}
Кроме того, при $A_{i-1}>0$
\begin{equation}\label{eq:coef-gap-log}
 \frac{a_i}{A_i}\le-\log\Bigl(1-\frac{a_i}{A_i}\Bigr)=\log\frac{A_i}{A_{i-1}} .
\end{equation}

\subsection{Невязка не больше \texorpdfstring{$4\varepsilon$}{4ε}: оценка по блокам}\label{sec:count1}

Первый способ подсчёта оценивает невязку после блока из $N$ последовательных итераций через невязку в начале блока.

\begin{lemma}[оценка по блокам]\label{lem:block}
Пусть $\norm{x^0-x^\ast}\le R$, алгоритм~\ref{alg:am-gap} запущен с $a_0=0$ и $r=R$, и итерации $k+1,\dots,k+N$, $N\ge1$, прошли без срабатывания критерия остановки. Тогда
\begin{equation}\label{eq:block}
 f(x^{k+N})-f^\ast\le\frac{96L_0R^2}{N^2}+\frac{3072\,L_1^2R^2}{N^3}\bigl(f(x^k)-f^\ast\bigr).
\end{equation}
\end{lemma}

\begin{proof}
Все суммы в доказательстве берутся по $i=k+1,\dots,k+N$.
\begin{itemize}
\item \emph{Невязка через веса.} Имеем
\begin{align*}
 f(x^{k+N})-f^\ast
 &\by{\le}{\eqref{eq:energy-gap}}\frac{R^2}{2A_{k+N}}
 \le\frac{R^2}{2\bigl(\sqrt{A_{k+N}}-\sqrt{A_k}\bigr)^2}\\
 &\by{\le}{\eqref{eq:coef-gap}}\frac{64R^2}{\Bigl(\sum_i\bigl(L_0+L_1\norm{\nabla f(y^i)}\bigr)^{-1/2}\Bigr)^2} .
\end{align*}
\item \emph{Неравенство Гёльдера.} С показателями $3/2$ и $3$
\begin{align*}
 N&=\sum_i\bigl(L_0+L_1\norm{\nabla f(y^i)}\bigr)^{-1/3}\bigl(L_0+L_1\norm{\nabla f(y^i)}\bigr)^{1/3}\\
 &\le\Bigl(\sum_i\bigl(L_0+L_1\norm{\nabla f(y^i)}\bigr)^{-1/2}\Bigr)^{2/3}\Bigl(\sum_i\bigl(L_0+L_1\norm{\nabla f(y^i)}\bigr)\Bigr)^{1/3} .
\end{align*}
Подставляя в предыдущую оценку, получаем
\begin{equation}\label{eq:block-g}
\begin{aligned}
 f(x^{k+N})-f^\ast
 &\le\frac{64R^2}{N^3}\sum_i\bigl(L_0+L_1\norm{\nabla f(y^i)}\bigr)\\
 &=\frac{64L_0R^2}{N^2}+\frac{64L_1R^2}{N^3}\sum_i\norm{\nabla f(y^i)} .
\end{aligned}
\end{equation}
\item \emph{Сумма градиентов мала, если мала невязка.} Суммируя неравенства~\eqref{eq:gap-monotone} по блоку, получаем
\[
 \sum_i\frac{\norm{\nabla f(y^i)}^2}{L_0+L_1\norm{\nabla f(y^i)}}\le8\bigl(f(x^k)-f^\ast\bigr).
\]
По неравенству КБШ
\begin{align*}
 \Bigl(\sum_i\norm{\nabla f(y^i)}\Bigr)^2
 &\le\sum_i\frac{\norm{\nabla f(y^i)}^2}{L_0+L_1\norm{\nabla f(y^i)}}\cdot\sum_i\bigl(L_0+L_1\norm{\nabla f(y^i)}\bigr)\\
 &\le8\bigl(f(x^k)-f^\ast\bigr)\Bigl(NL_0+L_1\sum_i\norm{\nabla f(y^i)}\Bigr).
\end{align*}
Если бы сумма $\sum_i\norm{\nabla f(y^i)}$ превышала $\sqrt{8NL_0(f(x^k)-f^\ast)}+8L_1(f(x^k)-f^\ast)$, то её квадрат превышал бы произведение этой величины на саму сумму, которое не меньше правой части последнего неравенства. Поэтому
\[
 \sum_i\norm{\nabla f(y^i)}\le\sqrt{8NL_0\bigl(f(x^k)-f^\ast\bigr)}+8L_1\bigl(f(x^k)-f^\ast\bigr).
\]
\item \emph{Подстановка.} Подставим эту оценку в~\eqref{eq:block-g}. Второе слагаемое в~\eqref{eq:block-g} станет не больше
\[
 \frac{64L_1R^2}{N^3}\sqrt{8NL_0\bigl(f(x^k)-f^\ast\bigr)}+\frac{512\,L_1^2R^2\bigl(f(x^k)-f^\ast\bigr)}{N^3} .
\]
По неравенству между средним арифметическим и средним геометрическим первое из этих слагаемых не больше
\[
 \frac{32L_0R^2}{N^2}+\frac{256\,L_1^2R^2\bigl(f(x^k)-f^\ast\bigr)}{N^3} .
\]
Поэтому
\[
 f(x^{k+N})-f^\ast\le\frac{96L_0R^2}{N^2}+\frac{768\,L_1^2R^2\bigl(f(x^k)-f^\ast\bigr)}{N^3} .
\]
Так как $768\le3072$, получаем~\eqref{eq:block}.\qedhere
\end{itemize}
\end{proof}

Оценка~\eqref{eq:block} объясняет, откуда берутся два слагаемых: первое --- обычная скорость ускоренного метода, второе --- цена нелинейности кривизны, пропорциональная невязке в начале блока. Если начальная невязка неизвестна и может быть экспоненциально большой, второе слагаемое бесполезно. Это препятствие обходится двумя способами: добавкой $\Phi$, если известен только радиус (раздел~\ref{sec:upper}), и положительным начальным весом, если известна оценка невязки (следующий подраздел).

\subsection{Известная оценка невязки: положительный начальный вес}\label{sec:count2}

Пусть $\norm{x^0-x^\ast}\le R$ и известно число $\Delta>\varepsilon$, для которого $f(x^0)-f^\ast\le\Delta$. Запустим алгоритм~\ref{alg:am-gap} с параметрами
\begin{equation}\label{eq:weight-params}
\begin{aligned}
 &a_0=\frac{R^2}{2\Delta},\\
 &r=\sqrt2R,\\
 &\delta=\varepsilon .
\end{aligned}
\end{equation}
Условие леммы~\ref{lem:energy-gap} с $R_0=R$ и $\Delta_0=\Delta$ выполнено, поскольку $R^2+2a_0\Delta=2R^2=r^2$.

\begin{lemma}[подсчёт при положительном начальном весе]\label{lem:count-weight}
Если $K$ итераций алгоритма~\ref{alg:am-gap} с параметрами~\eqref{eq:weight-params} прошли без срабатывания критерия остановки и $f(x^K)-f^\ast>\varepsilon$, то
\[
 K<16\sqrt{\frac{L_0R^2}\varepsilon}+21\,(L_1R)^{2/3}\log^{2/3}\frac{2\Delta}\varepsilon .
\]
\end{lemma}

\begin{proof}
Правая часть~\eqref{eq:energy-gap} не больше $R^2/2+R^2/2=R^2$. Так как $f(x^K)-f^\ast>\varepsilon$, отсюда
\begin{equation}\label{eq:count-weight}
\begin{aligned}
 A_K&<\frac{R^2}\varepsilon,\\
 \sum_{i=1}^Ka_i^2\norm{\nabla f(y^i)}^2&\le2R^2 .
\end{aligned}
\end{equation}
Разделим итерации на две группы: с малым градиентом, $L_1\norm{\nabla f(y^i)}\le L_0$, и с большим, $L_1\norm{\nabla f(y^i)}>L_0$.
\begin{itemize}
\item \emph{Итерации с малым градиентом.} Для них $L_0+L_1\norm{\nabla f(y^i)}\le2L_0$, и по~\eqref{eq:coef-gap} каждая такая итерация увеличивает $\sqrt{L_0A_i}$ не менее чем на $1/16$. Значит, их число не больше
\[
 16\sqrt{L_0A_K}\by{<}{\eqref{eq:count-weight}}16\sqrt{\frac{L_0R^2}\varepsilon} .
\]
\item \emph{Итерации с большим градиентом.} Для них $L_0+L_1\norm{\nabla f(y^i)}<2L_1\norm{\nabla f(y^i)}$, поэтому
\begin{align*}
 a_i\norm{\nabla f(y^i)}&\by{=}{\eqref{eq:coef-gap}}\frac{A_i\norm{\nabla f(y^i)}}{32\bigl(L_0+L_1\norm{\nabla f(y^i)}\bigr)a_i}\\
 &\bypad{>}{}{\eqref{eq:coef-gap}}\frac{A_i}{64L_1a_i} .
\end{align*}
Отсюда и из~\eqref{eq:count-weight}
\[
 \sum_{L_1\norm{\nabla f(y^i)}>L_0}\Bigl(\frac{A_i}{a_i}\Bigr)^2<(64L_1)^2\cdot2R^2=8192\,L_1^2R^2 .
\]
Так как $A_{i-1}\ge a_0>0$, по~\eqref{eq:coef-gap-log}
\begin{align*}
 \sum_{L_1\norm{\nabla f(y^i)}>L_0}\frac{a_i}{A_i}
 &\le\log\frac{A_K}{a_0}
 \by{<}{\eqref{eq:count-weight},\,\eqref{eq:weight-params}}\log\frac{2\Delta}\varepsilon .
\end{align*}
По неравенству Гёльдера с показателями $3$ и $3/2$ число итераций с большим градиентом равно
\begin{align*}
 \sum_{L_1\norm{\nabla f(y^i)}>L_0}\Bigl(\frac{A_i}{a_i}\Bigr)^{2/3}\Bigl(\frac{a_i}{A_i}\Bigr)^{2/3}
 &\le\Bigl(\sum\Bigl(\frac{A_i}{a_i}\Bigr)^2\Bigr)^{1/3}\Bigl(\sum\frac{a_i}{A_i}\Bigr)^{2/3}\\
 &<(8192\,L_1^2R^2)^{1/3}\log^{2/3}\frac{2\Delta}\varepsilon\\
 &<21\,(L_1R)^{2/3}\log^{2/3}\frac{2\Delta}\varepsilon .\qedhere
\end{align*}
\end{itemize}
\end{proof}

\subsection{Итоговые верхние оценки}\label{sec:gap-upper}

\begin{theorem}[верхние оценки при ограниченной начальной невязке]\label{thm:upper-gap}
Пусть $f\in\F_{L_0,L_1}^d$, $\norm{x^0-x^\ast}\le R$ и $\varepsilon>0$.
\begin{enumerate}
\item[(а)] Если $f(x^0)-f^\ast\le4\varepsilon$, то алгоритм~\ref{alg:am-gap} с $a_0=0$, $r=R$, $\delta=\varepsilon$,
\[
 m=\Bigl\lceil\log_2\Bigl(128\Bigl(1+\frac{L_0R^2}\varepsilon+L_1R\Bigr)^2\Bigr)\Bigr\rceil
\]
и числом итераций
\begin{align*}
 K&=\max\Bigl\{1,\Bigl\lceil\sqrt{\frac{192L_0R^2}\varepsilon}\,\Bigr\rceil,\bigl\lceil(24576\,L_1^2R^2)^{1/3}\bigr\rceil\Bigr\}\\
 &\le3+14\sqrt{\frac{L_0R^2}\varepsilon}+30(L_1R)^{2/3}
\end{align*}
выдаёт точку с невязкой по функции не больше $\varepsilon$, сделав не более $K(m+2)$ обращений к оракулу.
\item[(б)] Если известно число $\Delta>\varepsilon$ с $f(x^0)-f^\ast\le\Delta$, то алгоритм~\ref{alg:am-gap} с параметрами~\eqref{eq:weight-params},
\[
 m=\Bigl\lceil\log_2\Bigl(128\Bigl(1+\frac{2L_0R^2}\varepsilon+\sqrt2L_1R\Bigr)^2\Bigr)\Bigr\rceil
\]
и числом итераций
\[
 K=\Bigl\lceil16\sqrt{\frac{L_0R^2}\varepsilon}+21\,(L_1R)^{2/3}\log^{2/3}\frac{2\Delta}\varepsilon\Bigr\rceil
\]
выдаёт такую точку, сделав не более $K(m+2)$ обращений к оракулу. При $\Delta\le\varepsilon$ достаточно вернуть $x^0$.
\end{enumerate}
В обоих случаях метод использует только градиенты и не использует значение $f^\ast$. В случае~(а) он не использует и оценку начальной невязки.
\end{theorem}

\begin{proof}
\leavevmode
\begin{itemize}
\item \emph{Пункт~(а).} Если сработал критерий остановки, выход корректен по лемме~\ref{lem:search}. Иначе оценка~\eqref{eq:block} с $k=0$ и $f(x^0)-f^\ast\le4\varepsilon$ даёт
\[
 f(x^K)-f^\ast\le\frac{96L_0R^2}{K^2}+\frac{3072\,L_1^2R^2\cdot4\varepsilon}{K^3}\by{\le}{(\mathrm a)}\frac\varepsilon2+\frac\varepsilon2 ,
\]
где (а) --- выбор $K$. Оценка для $K$ следует из неравенств $\sqrt{192}<14$ и $24576^{1/3}<30$.
\item \emph{Пункт~(б).} При $\Delta\le\varepsilon$ возвращаем $x^0$. Иначе по лемме~\ref{lem:count-weight} после $K$ итераций либо сработал критерий остановки и выход корректен по лемме~\ref{lem:search}, либо $f(x^K)-f^\ast\le\varepsilon$.
\end{itemize}
В обоих пунктах число обращений не больше $K(m+2)$.
\end{proof}

Число шагов поиска $m+2$ растёт лишь логарифмически по $L_0R^2/\varepsilon$ и $L_1R$, поэтому теорема даёт для класса с начальной невязкой не больше $4\varepsilon$ и для класса с известной оценкой начальной невязки $\Delta$ соответственно оценки
\begin{align*}
 &O\Bigl(\Bigl[1+\sqrt{\frac{L_0R^2}\varepsilon}+(L_1R)^{2/3}\Bigr]\log\Bigl(2+\frac{L_0R^2}\varepsilon+L_1R\Bigr)\Bigr),\\
 &O\Bigl(\Bigl[1+\sqrt{\frac{L_0R^2}\varepsilon}+(L_1R)^{2/3}\log^{2/3}\frac{2\Delta}\varepsilon\Bigr]\log\Bigl(2+\frac{L_0R^2}\varepsilon+L_1R\Bigr)\Bigr).
\end{align*}
В случае~(б) перед запуском $\Delta$ можно заменить на минимум из $\Delta$ и правой части~\eqref{eq:radial-gap}. Вместе с теоремой~\ref{thm:lower-gap} первая оценка даёт при $L_0R^2\ge\varepsilon$
\[
 1+\mathcal N(\mathcal P_{4\varepsilon})=\widetilde\Theta\Bigl(1+\sqrt{\frac{L_0R^2}\varepsilon}+(L_1R)^{2/3}\Bigr).
\]
Совпадающая с~(б) нижняя оценка при $\Delta$ между $4\varepsilon$ и правой частью~\eqref{eq:radial-gap} неизвестна.
